% GENERATED FILE. Edit canonical lessons, then run npm run generate:books.
% canonical-source-hash: fnv1a64:c592e6b6a9afe046
% canonical-lessons: TE-C389-vislesincu, TE-C389-vislesana, TE-C389-parisilincu, TE-C389-polika, TE-C389-vyatyasam, TE-R389-review-382-389

\chapter{Analysing and Comparing}
\label{ch:te-analysing-and-comparing}
\hlchaptermodality{389}

\begin{chapteropening}
\textbf{By the end of this chapter:} \emph{I can analyse, examine and compare.}

Complete the last lesson of chapter 389: I can analyse, examine and compare.
\end{chapteropening}

\section[\texorpdfstring{\te{విశ్లేషించు} (viślēṣiñcu)}{విశ్లేషించు (viślēṣiñcu)}]{\te{విశ్లేషించు} (viślēṣiñcu) — to analyse}

\label{lesson:TE-C389-vislesincu}



\begin{quote}
Before the new one: say the Telugu for certainty, then the Telugu for uncertainty.
\end{quote}

\subsection*{You'll want to know: \te{విశ్లేషించు}}
\textbf{\te{విశ్లేషించు}} — \emph{viślēṣiñcu} — \textquotedblleft{}to analyse\textquotedblright{}.

\begin{quote}
\te{మేము} \te{పరీక్ష} \te{ఫలితాలను} \te{విశ్లేషించాము}.
\end{quote}

\begin{quote}
\emph{mēmu parīkṣa phalitālanu viślēṣiñcāmu.}
\end{quote}

\begin{quote}
We analysed the exam results.
\end{quote}

Past (I …): \te{విశ్లేషించాను}, \emph{viślēṣiñcānu}.

To take a question apart to find out how it works.

\begin{cousinweb}
\textbf{\te{విశ్లేషించు}} (\emph{viślēṣiñcu}) is built on Sanskrit \emph{viśleṣa}, \textquotedblleft{}loosening, separation, dissolution\textquotedblright{} (Monier-Williams), with the Telugu verb ending -\te{ఇంచు}.
\end{cousinweb}

\begin{grammarlens}[title={What you've built}]
One more word for giving reasons and opinions.
\end{grammarlens}

\subsection*{Your turn}
\begin{itemize}
\raggedright
  \item \emph{Say it:} \emph{viślēṣiñcu}
  \item \emph{Say it:} \emph{viślēṣiñcu}, once more
  \item \emph{Say it:} \emph{mēmu parīkṣa phalitālanu viślēṣiñcāmu.} — We analysed the exam results.
  \item \emph{Recall:} say the Telugu for certainty, then the Telugu for uncertainty, then say \emph{viślēṣiñcu} again
\end{itemize}

\begin{tcolorbox}[breakable,colback=teal!4,colframe=teal!35!black,title={Before you move on}]
What does \te{విశ్లేషించు} mean? (\textquotedblleft{}To analyse\textquotedblright{}.) Say it once more.
\end{tcolorbox}
\section[\texorpdfstring{\te{విశ్లేషణ} (viślēṣaṇa)}{విశ్లేషణ (viślēṣaṇa)}]{\te{విశ్లేషణ} (viślēṣaṇa) — an analysis}

\label{lesson:TE-C389-vislesana}



\begin{quote}
Before the new one: say the Telugu for uncertainty, then the Telugu for to analyse.
\end{quote}

\subsection*{You'll want to know: \te{విశ్లేషణ}}
\textbf{\te{విశ్లేషణ}} — \emph{viślēṣaṇa} — \textquotedblleft{}an analysis\textquotedblright{}.

\begin{quote}
\te{అతని} \te{విశ్లేషణ} \te{చాలా} \te{స్పష్టంగా} \te{ఉంది}.
\end{quote}

\begin{quote}
\emph{atani viślēṣaṇa cālā spaṣṭaṅgā undi.}
\end{quote}

\begin{quote}
His analysis is very clear.
\end{quote}

From the same word as \te{విశ్లేషించు}.

\begin{cousinweb}
\textbf{\te{విశ్లేషణ}} (\emph{viślēṣaṇa}) is borrowed from Sanskrit \emph{viśleṣaṇa}, \textquotedblleft{}separation, dissolution\textquotedblright{} (Monier-Williams); Telugu uses it for taking a question apart, an analysis.
\end{cousinweb}

\begin{grammarlens}[title={What you've built}]
One more word for giving reasons and opinions.
\end{grammarlens}

\subsection*{Your turn}
\begin{itemize}
\raggedright
  \item \emph{Say it:} \emph{viślēṣaṇa}
  \item \emph{Say it:} \emph{viślēṣaṇa}, once more
  \item \emph{Say it:} \emph{atani viślēṣaṇa cālā spaṣṭaṅgā undi.} — His analysis is very clear.
  \item \emph{Recall:} say the Telugu for uncertainty, then the Telugu for to analyse, then say \emph{viślēṣaṇa} again
\end{itemize}

\begin{tcolorbox}[breakable,colback=teal!4,colframe=teal!35!black,title={Before you move on}]
What does \te{విశ్లేషణ} mean? (\textquotedblleft{}An analysis\textquotedblright{}.) Say it once more.
\end{tcolorbox}
\section[\texorpdfstring{\te{పరిశీలించు} (pariśīliñcu)}{పరిశీలించు (pariśīliñcu)}]{\te{పరిశీలించు} (pariśīliñcu) — to examine closely; to look into}

\label{lesson:TE-C389-parisilincu}



\begin{quote}
Before the new one: say the Telugu for to analyse, then the Telugu for an analysis.
\end{quote}

\subsection*{You'll want to know: \te{పరిశీలించు}}
\textbf{\te{పరిశీలించు}} — \emph{pariśīliñcu} — \textquotedblleft{}to examine closely; to look into\textquotedblright{}.

\begin{quote}
\te{అధికారి} \te{నా} \te{దరఖాస్తును} \te{పరిశీలించాడు}.
\end{quote}

\begin{quote}
\emph{adhikāri nā darakhāstunu pariśīliñcāḍu.}
\end{quote}

\begin{quote}
The officer examined my application.
\end{quote}

Past (I …): \te{పరిశీలించాను}, \emph{pariśīliñcānu}.

Careful, official checking: \te{దరఖాస్తు} \te{పరిశీలించు}, to go through an application.

\begin{cousinweb}
\textbf{\te{పరిశీలించు}} (\emph{pariśīliñcu}) is built on Sanskrit \emph{pariśīlana}, \textquotedblleft{}contact; constant occupation, study\textquotedblright{} (Monier-Williams), with the Telugu verb ending -\te{ఇంచు}.
\end{cousinweb}

\begin{grammarlens}[title={What you've built}]
One more word for giving reasons and opinions.
\end{grammarlens}

\subsection*{Your turn}
\begin{itemize}
\raggedright
  \item \emph{Say it:} \emph{pariśīliñcu}
  \item \emph{Say it:} \emph{pariśīliñcu}, once more
  \item \emph{Say it:} \emph{adhikāri nā darakhāstunu pariśīliñcāḍu.} — The officer examined my application.
  \item \emph{Recall:} say the Telugu for to analyse, then the Telugu for an analysis, then say \emph{pariśīliñcu} again
\end{itemize}

\begin{tcolorbox}[breakable,colback=teal!4,colframe=teal!35!black,title={Before you move on}]
What does \te{పరిశీలించు} mean? (\textquotedblleft{}To examine closely; to look into\textquotedblright{}.) Say it once more.
\end{tcolorbox}
\section[\texorpdfstring{\te{పోలిక} (pōlika)}{పోలిక (pōlika)}]{\te{పోలిక} (pōlika) — a likeness; a resemblance; a comparison}

\label{lesson:TE-C389-polika}



\begin{quote}
Before the new one: say the Telugu for an analysis, then the Telugu for to examine closely.
\end{quote}

\subsection*{You'll want to know: \te{పోలిక}}
\textbf{\te{పోలిక}} — \emph{pōlika} — \textquotedblleft{}a likeness; a resemblance; a comparison\textquotedblright{}.

\begin{quote}
\te{అన్నకు} \te{తమ్ముడికి} \te{చాలా} \te{పోలిక} \te{ఉంది}.
\end{quote}

\begin{quote}
\emph{annaku tammuḍiki cālā pōlika undi.}
\end{quote}

\begin{quote}
The two brothers look very much alike.
\end{quote}

Plural: \te{పోలికలు}, \emph{pōlikalu}.

From the same root as \te{పోల్చు}, to compare, which you know.

\begin{cousinweb}
\textbf{\te{పోలిక}} (\emph{pōlika}) is native Dravidian: DEDR 4597 lists Telugu \emph{pōlika}, \textquotedblleft{}resemblance\textquotedblright{}, beside \emph{pōl(u)cu}, \textquotedblleft{}to compare\textquotedblright{}, and Tamil \emph{pōl}, \textquotedblleft{}to resemble\textquotedblright{}.
\end{cousinweb}

\begin{grammarlens}[title={What you've built}]
One more word for giving reasons and opinions.
\end{grammarlens}

\subsection*{Your turn}
\begin{itemize}
\raggedright
  \item \emph{Say it:} \emph{pōlika}
  \item \emph{Say it:} \emph{pōlika}, once more
  \item \emph{Say it:} \emph{annaku tammuḍiki cālā pōlika undi.} — The two brothers look very much alike.
  \item \emph{Recall:} say the Telugu for an analysis, then the Telugu for to examine closely, then say \emph{pōlika} again
\end{itemize}

\begin{tcolorbox}[breakable,colback=teal!4,colframe=teal!35!black,title={Before you move on}]
What does \te{పోలిక} mean? (\textquotedblleft{}A likeness; a resemblance; a comparison\textquotedblright{}.) Say it once more.
\end{tcolorbox}
\section[\texorpdfstring{\te{వ్యత్యాసం} (vyatyāsaṁ)}{వ్యత్యాసం (vyatyāsaṁ)}]{\te{వ్యత్యాసం} (vyatyāsaṁ) — a distinction; a contrast}

\label{lesson:TE-C389-vyatyasam}



\begin{quote}
Before the new one: say the Telugu for to examine closely, then the Telugu for a likeness.
\end{quote}

\subsection*{You'll want to know: \te{వ్యత్యాసం}}
\textbf{\te{వ్యత్యాసం}} — \emph{vyatyāsaṁ} — \textquotedblleft{}a distinction; a contrast\textquotedblright{}.

\begin{quote}
\te{రెండు} \te{ధరల} \te{మధ్య} \te{పెద్ద} \te{వ్యత్యాసం} \te{ఉంది}.
\end{quote}

\begin{quote}
\emph{reṇḍu dharala madhya pedda vyatyāsaṁ undi.}
\end{quote}

\begin{quote}
There is a big difference between the two prices.
\end{quote}

Formal beside \te{తేడా}, which you know.

\begin{cousinweb}
\textbf{\te{వ్యత్యాసం}} (\emph{vyatyāsaṁ}) is borrowed from Sanskrit \emph{vyatyāsa}, \textquotedblleft{}exchange; change, inverted order\textquotedblright{} (Monier-Williams). Telugu uses it for the difference between two things.
\end{cousinweb}

\begin{grammarlens}[title={What you've built}]
One more word for giving reasons and opinions.
\end{grammarlens}

\subsection*{Your turn}
\begin{itemize}
\raggedright
  \item \emph{Say it:} \emph{vyatyāsaṁ}
  \item \emph{Say it:} \emph{vyatyāsaṁ}, once more
  \item \emph{Say it:} \emph{reṇḍu dharala madhya pedda vyatyāsaṁ undi.} — There is a big difference between the two prices.
  \item \emph{Recall:} say the Telugu for to examine closely, then the Telugu for a likeness, then say \emph{vyatyāsaṁ} again
\end{itemize}

\begin{tcolorbox}[breakable,colback=teal!4,colframe=teal!35!black,title={Before you move on}]
What does \te{వ్యత్యాసం} mean? (\textquotedblleft{}A distinction; a contrast\textquotedblright{}.) Say it once more.
\end{tcolorbox}
\section[(dialogue)]{Reasons and opinions}

\label{lesson:TE-R389-review-382-389}



\begin{quote}
Say the Telugu for an idea and an outlook.
\end{quote}

\subsection*{Your turn}
Say these aloud:

\begin{itemize}
\raggedright
  \item an idea, an outlook, difference of opinion, to regard as, an angle
  \item a line of reasoning, evidence, to prove, logic, reasonable
  \item to suggest, to propose, a proposal, to recommend, a point
  \item to accept, consent, to agree with, to consent, backing
  \item to disagree, an objection, to object, opposition, to reject
  \item criticism, to criticise, to laud, to comment, to find fault with
  \item a misgiving, to have misgivings, to believe in, certainty, uncertainty
  \item to analyse, an analysis, to examine closely, a likeness, a distinction
\end{itemize}

\begin{tcolorbox}[breakable,colback=teal!4,colframe=teal!35!black,title={Before you move on}]
An idea? (\textbf{\te{ఆలోచన}}, \emph{ālōcana}.) An outlook? (\textbf{\te{దృక్పథం}}, \emph{dṛkpathaṁ}.) Difference of opinion? (\textbf{\te{అభిప్రాయభేదం}}, \emph{abhiprāyabhēdaṁ}.) To regard as? (\textbf{\te{భావించు}}, \emph{bhāviñcu}.) An angle? (\textbf{\te{కోణం}}, \emph{kōṇaṁ}.) A line of reasoning? (\textbf{\te{వాదన}}, \emph{vādana}.) Evidence? (\textbf{\te{సాక్ష్యం}}, \emph{sākṣyaṁ}.) To prove? (\textbf{\te{నిరూపించు}}, \emph{nirūpiñcu}.) Logic? (\textbf{\te{తర్కం}}, \emph{tarkaṁ}.) Reasonable? (\textbf{\te{సహేతుకమైన}}, \emph{sahētukamaina}.) To suggest? (\textbf{\te{సూచించు}}, \emph{sūciñcu}.) To propose? (\textbf{\te{ప్రతిపాదించు}}, \emph{pratipādiñcu}.) A proposal? (\textbf{\te{ప్రతిపాదన}}, \emph{pratipādana}.) To recommend? (\textbf{\te{సిఫార్సు} \te{చేయు}}, \emph{siphārsu cēyu}.) A point? (\textbf{\te{అంశం}}, \emph{aṁśaṁ}.) To accept? (\textbf{\te{అంగీకరించు}}, \emph{aṅgīkariñcu}.) Consent? (\textbf{\te{అంగీకారం}}, \emph{aṅgīkāraṁ}.) To agree with? (\textbf{\te{ఏకీభవించు}}, \emph{ēkībhaviñcu}.) To consent? (\textbf{\te{సమ్మతించు}}, \emph{sammatiñcu}.) Backing? (\textbf{\te{మద్దతు}}, \emph{maddatu}.) To disagree? (\textbf{\te{విభేదించు}}, \emph{vibhēdiñcu}.) An objection? (\textbf{\te{అభ్యంతరం}}, \emph{abhyantaraṁ}.) To object? (\textbf{\te{ఆక్షేపించు}}, \emph{ākṣēpiñcu}.) Opposition? (\textbf{\te{వ్యతిరేకత}}, \emph{vyatirēkata}.) To reject? (\textbf{\te{తిరస్కరించు}}, \emph{tiraskariñcu}.) Criticism? (\textbf{\te{విమర్శ}}, \emph{vimarśa}.) To criticise? (\textbf{\te{విమర్శించు}}, \emph{vimarśiñcu}.) To laud? (\textbf{\te{ప్రశంసించు}}, \emph{praśaṁsiñcu}.) To comment? (\textbf{\te{వ్యాఖ్యానించు}}, \emph{vyākhyāniñcu}.) To find fault with? (\textbf{\te{తప్పుబట్టు}}, \emph{tappubaṭṭu}.) A misgiving? (\textbf{\te{సందేహం}}, \emph{sandēhaṁ}.) To have misgivings? (\textbf{\te{సందేహించు}}, \emph{sandēhiñcu}.) To believe in? (\textbf{\te{విశ్వసించు}}, \emph{viśvasiñcu}.) Certainty? (\textbf{\te{నిశ్చయం}}, \emph{niścayaṁ}.) Uncertainty? (\textbf{\te{అనిశ్చితి}}, \emph{aniściti}.) To analyse? (\textbf{\te{విశ్లేషించు}}, \emph{viślēṣiñcu}.) An analysis? (\textbf{\te{విశ్లేషణ}}, \emph{viślēṣaṇa}.) To examine closely? (\textbf{\te{పరిశీలించు}}, \emph{pariśīliñcu}.) A likeness? (\textbf{\te{పోలిక}}, \emph{pōlika}.) A distinction? (\textbf{\te{వ్యత్యాసం}}, \emph{vyatyāsaṁ}.)
\end{tcolorbox}
